\documentclass{article}
\input{../packages.tex}
\begin{document}
    \begin{titlepage}
        \begin{center}

        Міністерство освіти і науки України
        
        НТУУ «Київський політехнічний інститут»
        
        Фізико-технічний інститут
        \vspace{3.3cm}
        
        {\textbf{Блокчейн та основи криптовалют}\\Конспект 0. Бекграунд. Геш-функції. ЦП.}

        \vspace{8.5cm}

        \begin{flushright}
            \textbf{Виконав:}\\Студент 4-го курсу\\групи ФІ-21\\Климентьєв Максим\\
            \textbf{Перевірив:}\\\text{\_\_\_\_\_\_\_\_\_\_\_\_\_\_\_\_\_\_}
        \end{flushright}

        \end{center}
    \end{titlepage}
    \newpage

    \pagenumbering{gobble}
    \tableofcontents
    \cleardoublepage
    \pagenumbering{arabic}
    \setcounter{page}{3}

    \newpage
    \section{Геш-функція}
        $ H: V^* \longrightarrow V_n $

        Зазвичай $ \perp $ --- бітовий вектор довжини 0

        Колізія --- коли геш-функція дає один результат для двох різних повідомлень. Повідомлення $M$ та $M_1$ мають однакове значення геш-функції $H(M) = H(M_1)$, отже мають однакове значення і цифрового підпису $ S(H(M)) = S(H(M_1)) $, отже виходить що можна підмінити ніби людина, яка написала повідомлення $M$, написала повідомлення $M_1$.

        Геш-функція, стійка до колізій --- коли неможливо для майже всіх повідомлень знайти два таких повідомлення.

        Основні властивості (вимоги) геш-функції:
        \begin{enumerate}
            \item Стійкість до знаходження першого прообразу --- Якщо знаємо $ y \in V_n $, то обчислювально неможливо знайти таке $x: H(x) = y$.
            \item Стійкість до знаходження колізії для заданого значення $x$ (другого прообразу) --- Маємо $x, H(x)$ --- обчислювально неможливо знайти таке $x_1: H(x) = H(x_1)$
            \item Стійкість до знаходження колізії --- Неможливо знайти таку пару $x_1, x_2: H(x_1) = H(x_2)$
        \end{enumerate}

        Крім цих вимог також є мнемонічні властивості:
        \begin{enumerate}
            \item Якщо змінити прийнаймі один біт входу, то в середньому зміниться половина біт виходу.
            \item Кожен біт входу впливає на кожен біт виходу.
        \end{enumerate}
        
        \textbf{Метод Днів Народження.} $2^{\frac{n}{2}}$ --- кількість хешувань для знаходження другого прообразу з ймовірністю суттєво вище 0.

        Після появи блокчейну з'явилась ще одна властивість: Складно знайти такий $ x $ для заданого $ v $, що $H(v||x) = y$

        Візьмемо $A \subset V_n$, $|A| << 2^n$ --- тоді обчислювально складно підібрати $H(V||x) \in A$

        У біткоїні приблизно $A$ вибирається: $A = \{ \underbrace{0 ... 0}_{70}  \underbrace{? ... ?}_{186} \}$. Для підбору $ x $ в середньому треба виконати $2^{70}$ спроб.

    \section{Цифровий підпис}
        $S: V_n \longrightarrow V_n$

        $S_k(H(M))$, де $k$ --- секретний ключ підписувача 
        
        Маємо Абонента
        $A: k, K$, де $K$ --- публічний ключ. 
        
        Для підписування треба лише $k$, для перевірки ж використовуєтся $K$.
        
        Для перевірки треба обчислити два різних перетворення та порівняти їх.
        
        $B: V_K(... =? ...) = 1$

        $Z_n = \{ 0, ..., n-1\}$

        $Z_p^*$ --- поле, $g$ --- генератор

        $a^{-1} \mod = b: ab \mod n = 1$
    
    \section{Мала теорема Ферма}
        $a^{p-1}\mod{p} = 1$

        \section{Задача Дискретного логарифмування (DLP)}
        $Z_p^*$, $g$ --- генератор

        $y \in Z_p^*$

        Сама задача: $? 0 <= k <= p - 2: g^k \mod{p} = y?$

        $k = ind(y, g, p)$

    \section{Алгоритм ЦП Ель-Гамаля}
        Перевірка публічного ключа.

        Обидва знають --- $Z_p^*$, $g$ --- генератор
        
        $A$ --- $a$: секретний ключ, $ 2 < a < p-2 $,
        
        \quad \quad $K = g^a\mod{p}$: відкритий ключ.
        
        $B$ --- $h$: значення надане сертифікатом, яке належить $A$

        $A$ хоче надіслати $M$, для побудови підпису $A$ обирає випадкове число $ r (2 < r < p-1), (r, p-1) = 1$ (взаємно просте).

        Далі $A$ обчислює величини $ s_1 = g^r \mod{p}$, $s_2 = (M - a s_1) r^{-1} \mod{(p-1)} $

        $A$ формує підпис $ S = (s_1, s_2) $ і надсилає $ B $ разом з повідомленням.

        Боб перевіряє чи $ g^M =? h^{s_1} s_1^{s_2} $

        Чому така формула: $ h^{s_1} s_1^{s_2} = g^{a s_1} g^{r s_2} = g^{M - a s_1} = g^M $
\end{document}
